\documentclass[11pt,a4paper]{article}

% ---------- Packages ----------
\usepackage[margin=1in]{geometry}
\usepackage{amsmath, amssymb, amsthm, mathtools}
\usepackage{microtype}
\usepackage{algorithm}
\usepackage{algpseudocode}
\usepackage[colorlinks=true,
            linkcolor=blue!50!black,
            urlcolor=blue!50!black,
            citecolor=blue!50!black]{hyperref}

% ---------- Theorem environments ----------
\theoremstyle{plain}
\newtheorem{theorem}{Theorem}[section]
\newtheorem{proposition}[theorem]{Proposition}
\newtheorem{lemma}[theorem]{Lemma}
\newtheorem{corollary}[theorem]{Corollary}

\theoremstyle{definition}
\newtheorem{definition}[theorem]{Definition}

\theoremstyle{remark}
\newtheorem{remark}[theorem]{Remark}
\newtheorem{example}[theorem]{Example}

% ---------- Math shortcuts ----------
\newcommand{\R}{\mathbb{R}}
\newcommand{\N}{\mathbb{N}}
\newcommand{\E}{\mathbb{E}}
\newcommand{\Var}{\operatorname{Var}}
\newcommand{\F}{\mathcal{F}}
\newcommand{\Prob}{\mathbb{P}}
\newcommand{\Q}{\mathbb{Q}}
\newcommand{\dd}{\mathop{}\!\mathrm{d}}
\newcommand{\eqdef}{\overset{\mathrm{def}}{=}}
\newcommand{\one}{\mathbf{1}}

% ---------- Document metadata ----------
\title{Monte Carlo Pricing of a European Call:\\
       Exact Simulation under Geometric Brownian Motion}
\author{Quant Finance Portfolio --- Phase 2, Block 1, Algorithm 1}
\date{\today}

\begin{document}
\maketitle

\begin{abstract}
We construct, analyse, and prepare for implementation the standard
Monte Carlo estimator of a European call option under geometric
Brownian motion, using exact simulation of the terminal asset price.
The setting is analytically degenerate: the price admits the
Black--Scholes closed form derived in Phase~1, and the algorithm is
therefore presented not for its practical necessity but as the
canonical Monte Carlo problem against which every more elaborate method
of Phase~2 will be measured. Three deliverables emerge. First, the
estimator is exhibited and the regularity hypotheses of the Block~0
framework are verified explicitly. Second, the variance of the
discounted payoff is computed in closed form, providing both an a
priori sample-size criterion and a sharp baseline for the variance
reductions of Block~2. Third, a triangulated validation strategy is
articulated, anchored on the Phase~1 Black--Scholes formula, on the
$O(n^{-1/2})$ asymptotic error rate, and on the agreement between
sample variance and its closed-form value.
\end{abstract}

\tableofcontents

\section{Introduction}
\label{sec:intro}

The European call option under geometric Brownian motion is the most
analytically tractable contingent claim in mathematical finance: its
price is given by the Black--Scholes formula
\begin{equation}
\label{eq:bs}
C_{\mathrm{BS}} \;=\; S_0\,N(d_1) \;-\; K e^{-rT}\,N(d_2),
\end{equation}
with the standard definitions of $d_1, d_2$ recalled in
Section~\ref{sec:setting}. Phase~1 of this project derived this formula
through the equivalent martingale measure framework and implemented it
in Python and C++.

The natural question is therefore why one would price this option by
Monte Carlo simulation at all. The honest answer is that, in this
setting, one would not---outside of pedagogical or validation contexts.
The algorithm developed in this writeup is the simplest possible
instance of the Monte Carlo machinery that will support every other
method in Phase~2, and its value lies in the role it plays as a
foundation for what follows.

That role is threefold. First, the algorithm is the canonical reference
implementation of a Monte Carlo \emph{engine}: a function that consumes
a model, a payoff, a sample size, and a random source, and returns a
point estimate together with a quantified uncertainty. The architecture
designed here---the partitioning of responsibilities between sampling,
payoff evaluation, and statistical reduction---will recur essentially
unchanged throughout the phase. Second, this is the only algorithm of
the phase entirely free of systematic error: the exact transition law
of the underlying is known and used directly, so no time-discretisation
bias enters, and no nested estimation introduces regression bias. Any
discrepancy between the estimator and the Black--Scholes price is
therefore unambiguously attributable to sampling error, which makes
this algorithm an unusually clean diagnostic for implementation
correctness. Third, because the distribution of the discounted payoff
is known explicitly, its variance admits a closed-form expression,
which we derive in Section~\ref{sec:variance}. This expression is the
denominator against which the variance reductions of Block~2 will be
measured, and it provides an additional independent test of the
implementation.

The remainder of the document is organised as follows.
Section~\ref{sec:setting} fixes notation. Section~\ref{sec:exact}
records the closed-form solution of the Black--Scholes SDE and
formalises the notion of exact simulation that the title invokes.
Section~\ref{sec:estimator} presents the estimator, verifies the
second-moment hypothesis required by the Block~0 framework, and
specialises the confidence interval to the present setting.
Section~\ref{sec:variance} derives the closed-form variance of the
discounted payoff and discusses its implications for sample-size
selection. Section~\ref{sec:validation} articulates the validation
strategy. Section~\ref{sec:algorithm} states the algorithm formally.
Section~\ref{sec:ahead} situates the algorithm within the rest of
Phase~2.

\section{Setting and notation}
\label{sec:setting}

\subsection{Risk-neutral dynamics}

We work on a filtered probability space
$(\Omega, \F, \{\F_t\}_{t \in [0,T]}, \Q)$ supporting a one-dimensional
$\Q$-Brownian motion $W$. Under $\Q$, the price process of the
underlying asset satisfies the stochastic differential equation
\begin{equation}
\label{eq:gbm-sde}
\dd S_t \;=\; r\, S_t \dd t \;+\; \sigma\, S_t \dd W_t,
\qquad S_0 > 0,
\end{equation}
where $r \in \R$ is the constant risk-free rate and $\sigma > 0$ is the
constant volatility. The numeraire is the money-market account
$B_t = e^{rt}$, so that the deflated price $S_t / B_t$ is a
$\Q$-martingale, and the discounted payoff of any contingent claim
admits a $\Q$-expectation interpretation.

\subsection{The European call payoff}

A European call with strike $K > 0$ and maturity $T > 0$ pays
$(S_T - K)^+$ at time $T$. Its time-zero price is the
discounted expectation
\begin{equation}
\label{eq:price}
C \;=\; \E^{\Q}\!\left[\, e^{-rT}\,(S_T - K)^+ \,\right]
\;=\; \E[\, Y \,],
\end{equation}
where we have set
\begin{equation}
\label{eq:Y-defn}
Y \;\eqdef\; e^{-rT}\,(S_T - K)^+
\end{equation}
and adopt the convention that, where unambiguous, the superscript on
$\E$ is dropped.

\subsection{The Black--Scholes constants}

We use the notation
\begin{equation}
\label{eq:d1d2}
d_1 \;\eqdef\;
\frac{\log(S_0/K) + \big(r + \tfrac{1}{2}\sigma^2\big)T}{\sigma\sqrt{T}},
\qquad
d_2 \;\eqdef\; d_1 - \sigma\sqrt{T}.
\end{equation}
The derivation of~\eqref{eq:bs} from~\eqref{eq:price} via the martingale
approach is given in the Phase~1 writeup and will not be repeated.

\section{Exact simulation of the terminal price}
\label{sec:exact}

\subsection{Closed-form solution of the SDE}

The SDE~\eqref{eq:gbm-sde} admits a closed-form solution,
\begin{equation}
\label{eq:gbm-solution}
S_t \;=\; S_0 \exp\!\left(
\big(r - \tfrac{1}{2}\sigma^2\big) t \;+\; \sigma\, W_t
\right),
\qquad t \in [0,T],
\end{equation}
which one verifies by applying It\^{o}'s lemma to
$f(t, x) = \log x$ and integrating. As a consequence,
$S_T$ is log-normally distributed: $\log S_T$ is Gaussian with mean
$\log S_0 + (r - \tfrac{1}{2}\sigma^2)T$ and variance $\sigma^2 T$.

\subsection{The single-normal sampling scheme}

Because $W_T \overset{d}{=} \sqrt{T}\, Z$ with $Z \sim N(0,1)$, exact
samples of $S_T$ are obtained by
\begin{equation}
\label{eq:sampling}
Z \sim N(0,1),
\qquad
S_T \;=\; S_0 \exp\!\left(
\big(r - \tfrac{1}{2}\sigma^2\big) T \;+\; \sigma\sqrt{T}\, Z
\right),
\end{equation}
each replication consuming exactly one standard normal draw. There is
no time-discretisation: the construction reaches $T$ in a single step.

\begin{remark}[The meaning of ``exact'']
\label{rem:exact-meaning}
The qualifier ``exact'' refers strictly to the absence of
time-discretisation error in the construction~\eqref{eq:sampling}: the
random variable on the right-hand side has \emph{exactly} the
distribution of $S_T$ under the model. It does not refer to the
estimator of the price, which carries the usual statistical error
inherent to Monte Carlo. The distinction matters because the algorithms
of Block~1.2 (Euler--Maruyama, Milstein) will not be exact in this
sense: their simulated terminal values will be drawn from a
distribution that approximates, but does not coincide with, the true
distribution of $S_T$.
\end{remark}

\section{The Monte Carlo estimator}
\label{sec:estimator}

\subsection{Definition and basic properties}

Let $Z_1, Z_2, \dots, Z_n$ be independent standard normal random
variables, and set
\begin{equation}
\label{eq:Si}
S_T^{(i)} \;=\; S_0 \exp\!\left(
\big(r - \tfrac{1}{2}\sigma^2\big) T \;+\; \sigma\sqrt{T}\, Z_i
\right),
\qquad
Y_i \;=\; e^{-rT}\,\big(S_T^{(i)} - K\big)^{+}.
\end{equation}
The Monte Carlo estimator of $C$ based on $n$ paths is
\begin{equation}
\label{eq:estimator}
\hat{C}_n \;\eqdef\; \frac{1}{n} \sum_{i=1}^{n} Y_i.
\end{equation}
By construction, $\hat{C}_n$ is unbiased: $\E[\hat{C}_n] = C$. By the
strong law of large numbers,
$\hat{C}_n \to C$ almost surely as $n \to \infty$.

\subsection{Verification of the second-moment hypothesis}
\label{sec:second-moment}

The asymptotic apparatus of Block~0---in particular the central limit
theorem and the resulting confidence interval---requires that
$\Var(Y) < \infty$, equivalently $\E[Y^2] < \infty$. We verify this
explicitly.

\begin{lemma}
\label{lem:second-moment}
The discounted payoff $Y$ defined in~\eqref{eq:Y-defn} satisfies
$\E[Y^2] < \infty$.
\end{lemma}

\begin{proof}
The pointwise inequality $0 \leq (S_T - K)^+ \leq S_T$ yields
$0 \leq Y \leq e^{-rT} S_T$. Hence
$\E[Y^2] \leq e^{-2rT}\, \E[S_T^2]$. From~\eqref{eq:gbm-solution},
\[
\E[S_T^2] \;=\; S_0^2\, \E\!\left[
\exp\!\big(2(r - \tfrac{1}{2}\sigma^2) T + 2\sigma\sqrt{T}\, Z\big)
\right]
\;=\; S_0^2\, \exp\!\big((2r + \sigma^2) T\big),
\]
where the last equality uses the moment-generating function of the
standard normal. Therefore
$\E[Y^2] \leq S_0^2\, e^{(\sigma^2 - 2r) T} \cdot e^{2rT}
= S_0^2\, e^{\sigma^2 T}$, which is finite.
\end{proof}

The hypotheses of the Block~0 framework are therefore satisfied without
restriction, and all conclusions of Sections~2.2--2.6 of that writeup
specialise directly to the present estimator.

\subsection{Confidence interval and stopping criterion}

Following the Block~0 prescription, the asymptotic
$(1-\beta)$-confidence interval for $C$ is
\begin{equation}
\label{eq:CI}
\hat{C}_n \;\pm\; z_{\beta/2}\,\frac{S_n}{\sqrt{n}},
\end{equation}
with $S_n^2 = (n-1)^{-1}\sum_{i=1}^n (Y_i - \hat{C}_n)^2$ the sample
variance of the payoffs and $z_{\beta/2}$ the upper $\beta/2$-quantile
of the standard normal. We denote the half-width by $h_n$. Two natural
stopping rules emerge:
\begin{enumerate}
\item Fix $n$ in advance (the more common choice in benchmarking work).
\item Continue sampling until $h_n$ falls below a prescribed
absolute or relative tolerance.
\end{enumerate}
Both will be supported by the implementation; the validation tests of
Section~\ref{sec:validation} use the first.

\section{Closed-form variance of the discounted payoff}
\label{sec:variance}

We now derive an explicit expression for $\Var(Y)$, which is finite by
Lemma~\ref{lem:second-moment} and computable in closed form because the
distribution of $S_T$ is known. The derivation reuses the
change-of-measure techniques familiar from the proof of~\eqref{eq:bs}
in Phase~1.

\subsection{Decomposition of the second moment}

On the event $\{S_T > K\}$, the payoff is $S_T - K$; elsewhere it
vanishes. Hence
\begin{equation}
\label{eq:Y2-expansion}
Y^2 \;=\; e^{-2rT}\, (S_T - K)^2\, \one_{\{S_T > K\}}
\;=\; e^{-2rT}\!\left[
S_T^{\,2}\, \one_{\{S_T > K\}}
- 2K\, S_T\, \one_{\{S_T > K\}}
+ K^2\, \one_{\{S_T > K\}}
\right].
\end{equation}
Three expectations must therefore be computed:
\begin{equation}
A \eqdef \E[\one_{\{S_T > K\}}],
\qquad
B \eqdef \E[S_T \one_{\{S_T > K\}}],
\qquad
D \eqdef \E[S_T^{\,2} \one_{\{S_T > K\}}].
\end{equation}

\subsection{The three expectations}

The first two are obtained as part of the standard derivation of the
Black--Scholes formula and we record them without re-deriving:
\begin{equation}
\label{eq:A-B}
A = N(d_2),
\qquad
B = S_0\, e^{rT}\, N(d_1).
\end{equation}
The third expectation is the new quantity. Writing $S_T$ in the
form~\eqref{eq:gbm-solution} with $W_T = \sqrt{T} Z$,
\[
S_T^{\,2}
\;=\; S_0^2 \exp\!\big(
(2r - \sigma^2) T + 2\sigma\sqrt{T}\, Z
\big),
\]
we have, with $\mu \eqdef (2r - \sigma^2) T$ and $\nu \eqdef
2\sigma\sqrt{T}$,
\begin{equation}
\label{eq:D-step1}
D \;=\; S_0^2\, \E\!\left[
e^{\mu + \nu Z}\, \one_{\{Z > -d_2\}}
\right],
\end{equation}
where the equivalence $\{S_T > K\} = \{Z > -d_2\}$ follows from
inverting~\eqref{eq:gbm-solution} for $Z$. Now perform the
exponential change of measure
\begin{equation}
\frac{\dd \tilde{\Q}}{\dd \Q} \;=\; e^{\nu Z - \nu^2/2}.
\end{equation}
Under $\tilde{\Q}$, the random variable $\tilde{Z} = Z - \nu$ is
standard normal (this is the elementary form of the Cameron--Martin
theorem). Substituting,
\[
\E\!\left[ e^{\nu Z}\, \one_{\{Z > -d_2\}} \right]
\;=\; e^{\nu^2/2}\, \tilde{\E}\!\left[ \one_{\{Z > -d_2\}} \right]
\;=\; e^{\nu^2/2}\, \tilde{\Prob}(\tilde{Z} > -d_2 - \nu)
\;=\; e^{\nu^2/2}\, N(d_2 + \nu).
\]
Therefore
\begin{equation}
\label{eq:D-result}
D \;=\; S_0^2\, e^{\mu + \nu^2/2}\, N(d_2 + \nu)
\;=\; S_0^2\, e^{(2r + \sigma^2) T}\, N\!\big(d_1 + \sigma\sqrt{T}\big),
\end{equation}
where we used $\mu + \nu^2/2 = (2r - \sigma^2)T + 2\sigma^2 T
= (2r + \sigma^2)T$ and
$d_2 + \nu = d_2 + 2\sigma\sqrt{T} = d_1 + \sigma\sqrt{T}$.

\subsection{The closed-form variance}

Substituting~\eqref{eq:A-B} and~\eqref{eq:D-result}
into~\eqref{eq:Y2-expansion} and combining with $\E[Y] = C_{\mathrm{BS}}$
yields the following result.

\begin{proposition}[Closed-form variance of the discounted call payoff]
\label{prop:closed-form-variance}
Let $Y$ be defined as in~\eqref{eq:Y-defn} under the dynamics
\eqref{eq:gbm-sde}. Then
\begin{equation}
\label{eq:closed-variance}
\Var(Y)
\;=\; S_0^{\,2}\, e^{\sigma^2 T}\, N\!\big(d_1 + \sigma\sqrt{T}\big)
\;-\; 2 K S_0\, e^{-rT}\, N(d_1)
\;+\; K^2\, e^{-2rT}\, N(d_2)
\;-\; C_{\mathrm{BS}}^{\,2},
\end{equation}
where $d_1, d_2$ are as in~\eqref{eq:d1d2} and $C_{\mathrm{BS}}$ as
in~\eqref{eq:bs}.
\end{proposition}

\begin{proof}
Combine~\eqref{eq:Y2-expansion}, \eqref{eq:A-B}, and~\eqref{eq:D-result}
to obtain $\E[Y^2]$:
\[
\E[Y^2] \;=\; e^{-2rT}\!\left[
S_0^2 e^{(2r + \sigma^2)T} N(d_1 + \sigma\sqrt{T})
- 2K \cdot S_0 e^{rT} N(d_1)
+ K^2 N(d_2)
\right].
\]
Distributing $e^{-2rT}$ gives the first three terms
of~\eqref{eq:closed-variance}. The identity
$\Var(Y) = \E[Y^2] - \E[Y]^2$ with $\E[Y] = C_{\mathrm{BS}}$ produces
the last term.
\end{proof}

\subsection{Implications for sample-size selection}

Proposition~\ref{prop:closed-form-variance} permits an a priori
calculation of the sample size required to achieve a prescribed
half-width $\varepsilon$ at confidence level $1-\beta$:
\begin{equation}
\label{eq:n-eps}
n_\varepsilon \;\approx\; \left(\frac{z_{\beta/2}}{\varepsilon}\right)^{\!2}
\Var(Y).
\end{equation}
This is one of the few cases in financial Monte Carlo where the
required sample size is known before any path is simulated. In every
later algorithm of the phase the variance must be estimated from a
preliminary run, and~\eqref{eq:n-eps} becomes a feedback loop rather
than a one-shot calculation.

The expression also serves a more lasting purpose: it is the
denominator $\Var_{\mathrm{vanilla}}$ of the variance reduction factors
\[
\mathrm{VRF} \;\eqdef\;
\frac{\Var_{\mathrm{vanilla}}}{\Var_{\mathrm{technique}}}
\]
that will be reported throughout Block~2 to quantify the gain of each
variance reduction technique.

\section{Validation strategy}
\label{sec:validation}

We follow the triangulation principle established in Phase~1: the
algorithm is subjected to several independent tests, each closing a
distinct potential failure mode. Passing any one of them is necessary
but not sufficient; passing all of them makes incorrectness highly
improbable, although strictly speaking unproven.

\paragraph{Test 1: containment of the Black--Scholes price.}
For a representative grid of parameters $(S_0, K, T, r, \sigma)$ and a
sample size $n$ moderate to large (e.g.\ $n = 10^4$ to $10^6$),
generate the estimate $\hat{C}_n$ together with its half-width $h_n$,
and verify that $C_{\mathrm{BS}}$ lies within
$[\hat{C}_n - h_n, \hat{C}_n + h_n]$. Repeating this experiment over
many independent seeds, the empirical containment frequency should
approach the nominal confidence level $1 - \beta$. This is the simplest
end-to-end check and tests the entire pipeline, from sampling to
statistical reduction.

\paragraph{Test 2: empirical convergence rate.}
For a fixed parameter configuration, evaluate $\hat{C}_n$ at sample
sizes $n_1 < n_2 < \cdots < n_K$ spanning several orders of magnitude.
Plot $\log |\hat{C}_n - C_{\mathrm{BS}}|$ against $\log n$. The
expected behaviour is a cloud of points lying around a line of slope
$-\tfrac{1}{2}$. To stabilise the plot against the variability of a
single run, average $|\hat{C}_n - C_{\mathrm{BS}}|^2$ over independent
seeds at each $n$ and plot
$\tfrac{1}{2}\log\!\big(\overline{|\hat{C}_n - C_{\mathrm{BS}}|^2}\big)$
against $\log n$, fitting a least-squares line to the resulting points.
The fitted slope should be close to $-\tfrac{1}{2}$.

\paragraph{Test 3: agreement of sample variance with closed form.}
For the same parameter configurations, compare the sample variance
$S_n^2$ produced by the simulation with the closed-form value
delivered by Proposition~\ref{prop:closed-form-variance}. The relative
error should decrease at the rate predicted by the variance of the
sample variance estimator, namely $O(n^{-1/2})$. A persistent
discrepancy at a level inconsistent with this rate indicates a bug in
either the sampling scheme or the closed-form computation; the test
discriminates between these two failure modes by comparison with a
third independent computation, performed at very large $n$.

\paragraph{Test 4: monotonicities and asymptotic limits.}
The Black--Scholes price is monotone non-decreasing in $S_0$, in
$\sigma$, and in $T$ (for fixed strike and rate); it is monotone
non-increasing in $K$. The Monte Carlo estimator must respect these
monotonicities up to statistical noise, which can be enforced sharply
by the use of common random numbers across parameter values (the same
seed produces the same sequence of $Z_i$ in both runs, and the
inequality holds path by path). The asymptotic limit
$\lim_{S_0 \to \infty} C / S_0 = 1$ should also be approximated, as
should $\lim_{S_0 \to 0} C = 0$.

\section{Algorithmic specification}
\label{sec:algorithm}

The algorithm is summarised below. The pseudocode is parametric in the
random source (a generator object accepting a seed) to enforce the
reproducibility discipline of Block~0, Section~2.5.

\begin{algorithm}[H]
\caption{Vanilla Monte Carlo for the European call (exact GBM)}
\label{alg:mc-vanilla}
\begin{algorithmic}[1]
\Require Spot $S_0$, strike $K$, maturity $T$, rate $r$, volatility
$\sigma$, sample size $n$, confidence level $1-\beta$, random
generator \texttt{rng} initialised with a recorded seed.
\Ensure Estimate $\hat{C}_n$, half-width $h_n$, sample variance
$S_n^2$.
\State $\mu \gets (r - \tfrac{1}{2}\sigma^2)\, T$ \Comment{drift of $\log S_T$}
\State $s \gets \sigma\, \sqrt{T}$ \Comment{stdev of $\log S_T$}
\State $D \gets e^{-rT}$ \Comment{discount factor}
\State Draw $Z_1, \dots, Z_n \sim N(0,1)$ independently using
\texttt{rng}.
\For{$i = 1$ to $n$}
  \State $S^{(i)} \gets S_0 \exp(\mu + s\, Z_i)$
  \State $Y_i \gets D \cdot \max(S^{(i)} - K,\, 0)$
\EndFor
\State $\hat{C}_n \gets \frac{1}{n} \sum_i Y_i$
\State $S_n^2 \gets \frac{1}{n-1} \sum_i (Y_i - \hat{C}_n)^2$
\State $h_n \gets z_{\beta/2}\, \sqrt{S_n^2 / n}$
\State \Return $(\hat{C}_n,\, h_n,\, S_n^2)$
\end{algorithmic}
\end{algorithm}

The implementation will compute steps~5--8 in vectorised form on the
entire vector $(Z_1, \dots, Z_n)$ at once, and the pseudocode is
written elementwise only for clarity. In neither Python nor C++ should
this loop appear as an explicit Python or C++ \texttt{for}: in Python
the entire computation is a sequence of NumPy operations, and in C++
it is a single contiguous loop with no allocation inside.

\section{Looking ahead}
\label{sec:ahead}

Three extensions of this algorithm structure the remainder of Phase~2.

\paragraph{Block~1.2: discretisation.} The exactness of the sampling
scheme~\eqref{eq:sampling} relies on the closed-form solution
\eqref{eq:gbm-solution} of the SDE, which is unavailable for almost
every model of practical interest beyond geometric Brownian motion. In
the next algorithm of the phase, the closed form is replaced by a
time-discretised approximation (Euler--Maruyama, then Milstein), and
systematic error of order $\Delta t$ or $\sqrt{\Delta t}$ enters
alongside the statistical error analysed here. The bias--variance
tradeoff of Section~5 of the Block~0 writeup becomes operational at
that point, and the optimal allocation of the computational budget
between sample size and time-step refinement is the central new
question.

\paragraph{Block~2: variance reduction.} Each technique introduced in
Block~2---antithetic variates, control variates, stratified sampling,
importance sampling---is implemented as a wrapper around the same
estimator structure developed here, and is benchmarked against the
present algorithm. The denominator of the variance reduction factor is
Proposition~\ref{prop:closed-form-variance}; the numerator is the
sample variance of the new estimator. The structural cleanliness of
the present algorithm---a single sampling step, a single payoff
evaluation, no per-replication state---makes it the natural reference
implementation in this comparison.

\paragraph{Block~4: greeks.} The pathwise and likelihood-ratio
estimators of price sensitivities reuse the sample paths
$\{S_T^{(i)}\}_{i=1}^n$ generated by the present algorithm, applying
to them a different functional in place of the discounted payoff. The
common-random-numbers technique used in conjunction with
finite-difference greeks requires that two parameter-perturbed
estimators consume the same sequence of $Z_i$, a requirement that the
seed-controlled sampling discipline of this algorithm makes immediate.

\section{Bibliographic notes}
\label{sec:biblio}

The exact-simulation Monte Carlo estimator for the European call under
geometric Brownian motion is treated, as a textbook starting example,
in Glasserman~\cite[Section~1.1.2 and Section~3.2]{Glasserman}. The
closed-form variance computation of
Proposition~\ref{prop:closed-form-variance} is elementary and appears
implicitly in many sources; an explicit derivation in the same spirit
as the present one can be found in Jäckel~\cite[Chapter~3]{Jackel}.
The variance reduction factor as the figure of merit for Monte Carlo
techniques is discussed in
Glasserman~\cite[Section~4.1]{Glasserman}.

\begin{thebibliography}{99}

\bibitem{Glasserman}
P.\ Glasserman,
\emph{Monte Carlo Methods in Financial Engineering},
Springer, 2004.

\bibitem{Jackel}
P.\ J\"ackel,
\emph{Monte Carlo Methods in Finance},
Wiley, 2002.

\bibitem{Etheridge}
A.\ Etheridge,
\emph{A Course in Financial Calculus},
Cambridge University Press, 2002.

\end{thebibliography}

\end{document}
